El estudio del tiempo de acceso y del rendimiento de las tecnologías de memoria,
resulta conveniente porque ayuda a comprender factores como la latencia, el 
ancho de banda y la organización de la memoria. Este análisis es importante
ya que ante el crecimiento de aplicaciones que requieren procesar grandes 
cantidades de datos, lo cual depende de una comunicación eficiente entre el 
procesador y la memoria.

Desde una perspectiva más práctica, la investigación permite identificar caracterśticas 
y limitaciones de tecnologías como DDR4, DDR5, y memorias de alto ancho de banda, facilitando
la compresión de los criterios de intervienen en la optimización del sistema de memoria.
Su valor teórico consiste en integrar resultados recientes sobre la relación entre latencia,
ancho de banda y rendimiento. Además, aunque no se pretender desarrollar un nuevo instrumento 
de medición, la revisión y comparación proporciona una metodología organizada para analizar
el comportamiento de distintas tecnologías.